\documentclass[preview]{standalone}

\usepackage{amsmath}
\usepackage{amssymb}
\usepackage{stellar}
\usepackage{definitions}

\begin{document}

\id{lo-strutturalismo}
\genpage

\section{Lo strutturalismo}

\begin{snippetdefinition}{strutturalismo-definition}{Strutturalismo}
    Lo \emph{strutturalismo} è un approccio psicologico fondato
    sull'elementarismo e sul metodo dell'introspezione. Compare nel 1898 in un
    articolo dello psicologo inglese \emph{E. B. Titchener}, allievo di Wundt.
    Il suo obiettivo è comprendere la struttura della mente scomponendo
    l'esperienza cosciente nei suoi elementi primari.
\end{snippetdefinition}

\begin{snippet}{strutturalismo-componenti-esperienza}
    Per Titchener, l'esperienza cosciente è costituita da percezioni, idee,
    emozioni o sentimenti. A queste corrispondono tre componenti fondamentali:
    \begin{itemize}
        \item le \emph{sensazioni}, alla base delle percezioni;
        \item le \emph{immagini mentali}, alla base delle idee;
        \item gli \emph{stati affettivi}, alla base delle emozioni.
    \end{itemize}
    Alle caratteristiche di qualità e intensità già individuate da Wundt,
    Titchener aggiunge durata e chiarezza.
\end{snippet}

\begin{snippetdefinition}{errore-dello-stimolo-definition}{Errore dello stimolo}
    L'\emph{errore dello stimolo} consiste nell'attribuire significati e
    valori soggettivi ai dati oggettivi dell'esperienza, invece di descrivere
    l'esperienza cosciente nei suoi elementi immediati.
\end{snippetdefinition}

\begin{snippet}{critiche-strutturalismo}
    Lo strutturalismo è stato criticato principalmente perché:
    \begin{itemize}
        \item è \emph{riduzionista}, poiché riconduce la complessità
        dell'esperienza umana a componenti elementari;
        \item è \emph{elementarista}, poiché non studia direttamente la totalità
        di un comportamento;
        \item è \emph{mentalista}, poiché analizza soprattutto resoconti verbali
        della coscienza e trascura soggetti che non possono produrli, come
        bambini o persone con gravi disturbi.
    \end{itemize}
\end{snippet}

\end{document}
